% Laget automatisk av elevark/lag_elevark.py fra sim/generator.html. Endringer her blir overskrevet.
\documentclass[11pt]{article}
\usepackage{elevark}
\begin{document}
\elevarkhode{Fysikk 2 · Elektromagnetisk induksjon}{Generator}{https://fysikklab.no/sim/generator.html}{Simuleringen}

Nesten all elektrisk energi lages i generatorer: En spole roterer i et magnetfelt, enten den drives av vann, vind eller damp. Fluksen gjennom spolen endrer seg hele tiden, og det blir indusert en vekselspenning. Hva bestemmer hvor stor spenningen blir?

\section*{Oppgaver}

\begin{oppgaveliste}
\oppgave \fase{gjett}{Gjett}I hvilken stilling er spenningen størst: når spoleplanet står vinkelrett på feltet, eller når det er parallelt med feltet?
\svar
\oppgave \fase{test}{Test}Trykk Start og sammenlign spolen med grafen. Hvor er fluksen størst, og hvor er spenningen størst? Stemmer det du gjettet?
\svar
\oppgave \fase{forklar}{Forklar}Vis at $ε = NBAω \sin (ωt)$ når $NΦ = NBA \cos (ωt)$. Hvorfor blir spenningen størst når fluksen er null?
\svar
\oppgave \fase{test}{Undersøk}Hva skjer med toppspenningen og med perioden når du dobler $f$? Når du dobler $N$, $B$ eller $A$?
\svar
\oppgave \fase{test}{Regn}Et strømnett har 50~Hz og effektivverdi 230~V. Hvor stor må NBA være? Finn en kombinasjon i simuleringen som gir omtrent dette.
\svar
\oppgave \fase{forklar}{Forklar}Kryss av for «Kommutator». Hva gjør den med spenningen? Hvorfor trenger en sykkeldynamo eller et kraftverk ikke kommutator?
\svar
\oppgave \fase{utfordring}{Utfordring}Når lyspæra er koblet til, går det strøm i spolen. Hvorfor blir det tyngre å dreie generatoren da? Hva betyr det for et vannkraftverk når mange slår på ovnen samtidig?
\svar
\end{oppgaveliste}

\end{document}
